\begin{frame}{이번 학기의 실습 환경: Gymnasium}
  \begin{columns}[T]
    \begin{column}{0.5\textwidth}
      \begin{itemize}
        \item \kb{Gymnasium} --- 옛 OpenAI Gym(Brockman et al., 2016)의
          후신. \textbf{표준} 인터페이스 라이브러리.
        \item 어떤 환경이든: 현재 상태를 관측(\texttt{reset})하고,
          행동을 하나 취하면(\texttt{step})
          \textbf{다음 상태$\cdot$보상$\cdot$종료 여부}가 돌아옴.
        \item 인터페이스가 통일되므로 2$\sim$11주차의 어떤 알고리즘이든
          \textbf{환경만 바꿔 끼우면} 재사용.
        \item 13$\sim$14주차의 MuJoCo$\cdot$Isaac Sim 로봇 환경도
          이 \texttt{reset}/\texttt{step} 패턴을 그대로 따름.
      \end{itemize}
    \end{column}
    \begin{column}{0.46\textwidth}
      \centering
      \includegraphics[width=\linewidth]{ref_gym.png}
      \vskip0.4em
      {\footnotesize Gym에 포함된 환경들 (CarRacing, Go, AirRaid,
      LunarLander, MountainCar, Walker2d, Doom, Reacher) --- 원 논문
      Figure 1.}
    \end{column}
  \end{columns}
\end{frame}
